\SourceLeaf{36}

\Needspace{8\baselineskip}
\begin{center}\small
BEMERKUNGEN ZU DEM AUFSATZE »BEWEIS DER IRRATIONALITÄT DER TANGENTEN
RATIONALER BÖGEN IN EINER NEUEN GESTALT«.
\end{center}
\par\noindent\hspace*{\parindent}Der erste einwandfreie Beweis, dass $\operatorname{tg}\dfrac{m}{n}$ für ganze Zahlen
$m, n$ irrational ist, wurde von \textsc{Legendre} in der vierten Note seiner »Éléments
de géométrie« (1794) entwickelt. Die \textsc{Gauss}'schen Ausführungen weichen in den
Grundgedanken nicht von dem \textsc{Legendre}schen Beweise ab. Neu ist die Art, wie
\textsc{Gauss} direct aus den Potenzreihen der $P_\theta$ die unbegrenzte Abnahme dieser
Grössen bei wachsendem $\theta$ nachweist. Über die hier in Betracht kommenden
Untersuchungen \textsc{Lamberts} und \textsc{Legendres}, sowie namentlich über die vom
letzteren a.\,a.\,O. gegebene Darstellung hat übrigens \textsc{Gauss} eine kritische
Besprechung aufgezeichnet, die jedoch hier nicht abgedruckt ist.

Eine andere Entwicklung, welche gleichfalls nicht abgedruckt ist, leitet aus den
Recursionsformeln der $P_\theta$ angenäherte Darstellungen von $\operatorname{tg} m$
durch Quotienten ganzer rationaler Functionen von $m$ ab. \textsc{Gauss} benutzt diese
Darstellungen zur Berechnung von Näherungswerthen von $\tfrac{1}{2}\pi$.

\begin{flushright}
\textsc{Fricke}.
\end{flushright}
